%% ============================================================
%% JURIX 2025 — Hybrid Sparse-Dense Retrieval for EU Climate Law
%%
%% DRAFT — uses standard LaTeX `article` class so it compiles on
%% any TeX installation (Overleaf, local, CI).
%%
%% FOR FINAL SUBMISSION: replace \documentclass{article} with
%% \documentclass{iosart2x} and restore the IOS Press frontmatter
%% macros (\inits, \fnm, \snm, \address, keyword/\sep, frontmatter
%% environment).  The iosart2x.cls is supplied by IOS Press with the
%% author kit.
%%
%% STANDALONE: bibliography is inline — no .bib file needed.
%% Entries marked [VERIFY] require author/venue confirmation.
%%
%% TODO (before submission):
%%   1. Add repository URL in §10.4.
%%   2. Verify all [VERIFY]-flagged bibliography entries.
%%   3. Restore IOS Press author-kit macros for camera-ready.
%% ============================================================
\documentclass[11pt,a4paper]{article}

\usepackage[top=2.5cm,bottom=2.5cm,left=2.5cm,right=2.5cm]{geometry}
\usepackage[T1]{fontenc}
\usepackage[utf8]{inputenc}
\usepackage{amsmath,amssymb}
\usepackage{booktabs}
\usepackage{multirow}
\usepackage{graphicx}
\usepackage{microtype}
\usepackage[ruled,vlined]{algorithm2e}
\usepackage[hidelinks]{hyperref}
\usepackage{xspace}
\usepackage{authblk}

%% ── Convenience macros ───────────────────────────────────────────
\newcommand{\HR}[1]{\mathrm{HR}@{#1}}
\newcommand{\MRR}[1]{\mathrm{MRR}@{#1}}
\newcommand{\nomic}{\texttt{nomic-embed-text-v1.5}\xspace}
\newcommand{\nomicshort}{nomic-embed\xspace}
\newcommand{\rrf}{RRF\xspace}

%% ── Title / authors / affiliations ───────────────────────────────
\title{Hybrid Sparse-Dense Retrieval for EU Climate Law:\\
       Evaluating BM25, \nomic and Reciprocal Rank Fusion\\
       on a Statute-Level Benchmark}

\author[1]{Rameez Wani}
\author[1]{Ananya Warior}
\affil[1]{School of Computing, Dublin City University, Dublin, Ireland\\
  \texttt{rameezshafat.wani2@mail.dcu.ie;\
          ananyasalil.warior2@mail.dcu.ie}}

\date{}

%% ============================================================
\begin{document}
\maketitle

\begin{abstract}
Grounded responses in legal AI applications depend critically on
first-stage retrieval: a system that fails to surface the correct
statutory provision cannot produce a verifiable, citation-supported
answer.  We evaluate three retrieval configurations---BM25 lexical
search, \nomic dense retrieval (768\,dim, 8{,}192-token context), and
a hybrid combining both via weighted Reciprocal Rank Fusion (\rrf,
$k{=}20$, dense weight~5)---on a manually constructed benchmark of
71~queries over EU Climate Law.  The corpus comprises 1{,}166
article-level provisions from 72~EU instruments extracted from the EU
CELLAR database via SPARQL and Formex~XML.  On the full 71-query
benchmark, the hybrid achieves $\HR{5}{=}0.972$ compared with
$\HR{5}{=}0.704$ for BM25 and $\HR{5}{=}0.958$ for \nomicshort alone;
the improvement over BM25 is significant (McNemar $p{<}0.001$) while
the marginal gain over \nomicshort alone is not ($p{=}0.317$).  Only
one query ($1.4\,\%$) is a hard failure for all three systems.  To the
best of our knowledge, this is the first systematic comparative study
of hybrid sparse-dense retrieval on EU statutory text, and the first
statute-level retrieval benchmark for EU Climate Law.  The findings
indicate that \nomicshort dense retrieval is already near-optimal for
this domain, and that lexical-semantic fusion provides a small but
consistent safety margin.
\end{abstract}

\medskip
\noindent\textbf{Keywords:} EU Climate Law;\enspace legal information
retrieval;\enspace hybrid retrieval;\enspace BM25;\enspace
nomic-embed;\enspace reciprocal rank fusion;\enspace
retrieval-augmented generation;\enspace statute-level benchmark.

% ============================================================
\section{Introduction}
\label{sec:intro}

As a motivating hypothetical, consider a legal analyst querying a
BM25-only system with \textit{``What monitoring obligations apply to
aircraft operators under the EU Emissions Trading System?''}\ A
BM25-only system might return top-ranked provisions from
Regulation~(EU)~2015/757 on maritime emissions monitoring
(\textsc{celex}~\texttt{32015R0757}) because the terms
\textit{monitoring}, \textit{operators}, and \textit{obligations}
achieve high BM25 scores against that instrument's operative
articles.  The correct provisions---Articles~14--15 of
Directive~2003/87/EC (\texttt{32003L0087})---could rank below
position~10 owing to aviation-specific vocabulary
(\textit{tonne-kilometre data}, \textit{verified emissions report})
that is absent from the query.  The downstream language model,
receiving context that contains no mention of aviation or
Article~14, would be unable to produce a grounded answer and would
either hallucinate a response or correctly invoke an
insufficient-context guard.

This hypothetical scenario illustrates why first-stage retrieval
quality is the rate-limiting factor for grounded legal AI
applications.  Retrieval-augmented generation (RAG) pipelines for
statute interpretation require that the retriever surface not merely
topically similar text, but the legally binding provisions that
\emph{answer} the query.  In EU Climate Law, this requirement is
compounded by legislative density (72 enacted instruments),
cross-instrument definitional dependencies (e.g., the CBAM
Regulation cross-references ETS allowance concepts defined in
Directive~2003/87/EC), highly technical domain-specific vocabulary
(\textit{free allocation}, \textit{carbon leakage},
\textit{corrigendum}), and the urgent policy context that makes
error-free source attribution a professional and legal necessity.

Sparse lexical retrieval (BM25) captures exact statutory terminology
but is susceptible to vocabulary mismatch~\cite{robertson2009bm25}.
Dense semantic retrieval (\nomicshort) generalises beyond exact matches but
may conflate legally distinct concepts sharing surface
similarity~\cite{nussbaum2024nomic}.  Hybrid fusion via Reciprocal Rank
Fusion (\rrf) combines both lists without requiring score
normalisation or labelled training data~\cite{cormack2009rrf},
making it particularly suitable for a domain where annotated
retrieval pairs are scarce.

\paragraph{Research Questions.}

\begin{description}
  \item[\textbf{RQ1.}] Does hybrid sparse-dense retrieval
    (BM25\,+\,\nomicshort with \rrf $k{=}20$, dense weight~5)
    improve the effectiveness and reliability of first-stage
    grounding for EU Climate Law compared with BM25-only and
    \nomicshort-only retrieval?
  \item[\textbf{RQ2.}] Is first-stage hybrid retrieval sufficient to
    support accurate and domain-relevant grounded responses in EU
    Climate Law, given retrieval coverage, hard-query failure rates,
    and efficiency overhead?
\end{description}

\paragraph{Contributions.}  This paper makes four contributions:

\begin{enumerate}
  \item A 71-query manually verified benchmark for EU Climate Law
    statute-level retrieval with binary CELEX-level relevance
    judgements (\S\ref{sec:eval-bench}).
  \item A reproducible article-level corpus of 1{,}166 provisions
    from 72 EU climate instruments, extracted from the EU CELLAR
    SPARQL endpoint via a documented two-phase pipeline
    (\S\ref{sec:corpus}).
  \item A three-way comparative evaluation of BM25, \nomicshort, and
    BM25{+}\nomicshort \rrf across effectiveness, reliability, and
    efficiency dimensions (\S\ref{sec:results}).
  \item A characterisation of first-stage retrieval failure modes
    and their implications for reliable source attribution in
    domain-specific LLM deployments (\S\ref{sec:error}).
\end{enumerate}

% ============================================================
\section{Related Work}
\label{sec:related}

\subsection{Lexical Retrieval in Legal IR}
\label{sec:related-lexical}

BM25 remains the dominant first-stage lexical retrieval model owing
to its probabilistic grounding and robust performance without
training data~\cite{robertson2009bm25}.  In legal retrieval, exact
match over statutory vocabulary is particularly valuable: defined
terms such as \textit{installation}, \textit{allowance}, and
\textit{tonne-kilometre} carry precise legal meanings that
synonymy-based generalisation may distort.  Mori and
Kano~\cite{mori2023} examine statute-level retrieval for legal
question answering and find that lexical signals remain strong for
definitional and obligation-type queries where the relevant provision
shares key vocabulary with the question.  The NOWJ team at
COLIEE~\cite{nowjcoliee2023} demonstrate that well-tuned BM25
baselines are competitive on statutory entailment tasks, establishing
a high bar for learned alternatives.  The fundamental limitation of
pure lexical retrieval---term mismatch between user query vocabulary
and statutory vocabulary---motivates the dense component evaluated in
this paper.

\subsection{Dense Retrieval and Legal Embeddings}
\label{sec:related-dense}

Dense retrieval encodes both queries and documents into a shared
embedding space, enabling semantic matching beyond exact term
overlap~\cite{nussbaum2024nomic}.  The \nomic model achieves strong
performance on long-context retrieval benchmarks through a
Matryoshka-style training objective that supports variable-dimension
embeddings~\cite{nussbaum2024nomic}; critically for this work, its
8{,}192-token context window accommodates long legislative articles
without arbitrary truncation for the majority of the corpus.  Rayo
et al.~\cite{rayo2024} evaluate dense retrieval for regulatory
compliance queries and identify distributional shift between
general-domain pre-training corpora and EU statutory text as a
factor limiting performance.  LegalRAG~\cite{legalrag2024} integrates
dense retrieval into a legal QA pipeline and reports that semantic
matching reduces hard failures on queries whose vocabulary diverges
from the document collection.  A shared limitation is that dense
models may retrieve thematically related but legally distinct
instruments---a failure mode characterised in \S\ref{sec:error}.

\subsection{Sparse Learned Representations}
\label{sec:related-splade}

SPLADE~\cite{formal2021splade} learns sparse representations by
applying a log-saturation activation over BERT vocabulary logits,
producing inverted-index-compatible vectors that encode both lexical
and expanded signals.  The SPLADE-v2 extension~\cite{formal2022splade2}
improves effectiveness through distillation from a dense teacher and
hard-negative sampling, substantially narrowing the gap between
sparse learned models and dense retrievers on standard BEIR
benchmarks.  However, both variants require domain-specific
fine-tuning to be effective: without in-domain annotated pairs,
SPLADE models exhibit vocabulary hallucination---expanding queries
towards general-domain concepts absent from the target corpus.  No
publicly available SPLADE model trained on EU statutory text is known
to the authors.  This gap justifies using terminology-preserving BM25
as the lexical component, and flags SPLADE fine-tuning on EU law as a
productive direction for future work.

\subsection{Hybrid Retrieval and Rank Fusion}
\label{sec:related-hybrid}

Hybrid retrieval combining sparse and dense ranked lists via
Reciprocal Rank Fusion was shown by Cormack et
al.~\cite{cormack2009rrf} to outperform individual rank learning
methods and Condorcet fusion on TREC collections without requiring
score normalisation.  HyPA-RAG~\cite{hyparag2024} applies hybrid
parameter-adaptive retrieval to AI legal and policy applications.
LRAGE~\cite{lrage2024} evaluates multiple hybrid retrieval
configurations in a legal RAG benchmark framework.  A recurring
theoretical motivation across these systems is that the independence
assumption underlying \rrf---that ranking errors of each retriever
are partially uncorrelated---may hold in legal IR because lexical and
semantic failure modes are empirically distinct; this is a testable
claim that \S\ref{sec:eval-reliability} evaluates.

\subsection{Legal RAG Systems}
\label{sec:related-legalrag}

Several systems address the end-to-end legal RAG pipeline.
CBR-RAG~\cite{cbrrag2023} augments retrieval with case-based
reasoning, using structural similarity between retrieved precedents
and the current query as a re-ranking signal.  LexRAG~\cite{lexrag2024}
combines lexical retrieval with generation for legal reasoning tasks
and reports that lexical precision is a stronger predictor of answer
correctness than semantic similarity on statutory interpretation
queries.  LexDrafter~\cite{lexdrafter2024} applies RAG to legislative
drafting assistance, finding that article-level chunking aligns with
drafting practice better than fixed-token windows.  Branting et
al.~\cite{legalstructurerag2024} exploit legal document structure
(recitals, chapters, articles, sub-paragraphs) as a retrieval signal,
reporting gains on tasks where the correct provision is embedded in a
specific structural position.  A consistent gap across this body of
work is the absence of systematic evaluation on EU statutory text;
most systems focus on common-law jurisdictions (US, UK) or case law
retrieval, with EU civil-law legislation---which has distinct
structural and terminological properties---largely unexplored.

\subsection{Retrieval Reliability and Source Attribution}
\label{sec:related-reliability}

Ravichander et al.~\cite{reliableretrievallegal2024} identify
retrieval reliability---the consistency with which a system surfaces
the correct provision across paraphrased variants of the same
query---as a distinct evaluation dimension from single-query
effectiveness, and argue that unreliable retrieval is a principal
driver of hallucination in grounded legal AI.  Zheng et
al.~\cite{reasoninglegal2024} frame legal retrieval as a reasoning
problem in which the first-stage retriever must resolve implicit
definitional dependencies between instruments before a
chain-of-thought generator can be applied.  Both works underscore
that for legally consequential applications---where every claim in an
answer must be attributed to a specific, verifiable statutory
provision---single-point IR metrics understate the practical risk
introduced by retrievers with uneven per-query coverage.  We address
this gap directly through the reliability dimension in
\S\ref{sec:eval-reliability}.

\subsection{Evaluation Frameworks for Legal Retrieval}
\label{sec:related-eval}

The COLIEE benchmark~\cite{nowjcoliee2023} is the most widely used
legal retrieval evaluation, covering statute retrieval and legal
entailment tasks primarily over the Japanese Civil Code and Canadian
case law.  LRAGE~\cite{lrage2024} proposes a framework for evaluating
legal RAG systems that decomposes performance into retrieval coverage,
attribution accuracy, and answer faithfulness.  HyPA-RAG~\cite{hyparag2024}
includes an evaluation component for legal and policy question
answering.  None of these frameworks provides a benchmark for EU
Climate Law statute retrieval, where the target collection spans
72~instruments in the CELLAR repository and queries must resolve
cross-instrument definitional dependencies.  The benchmark
constructed in this paper fills this gap and is designed to be
directly reproducible from the publicly available CELLAR SPARQL
endpoint.

% ============================================================
\section{Research Gap and Contributions}
\label{sec:gap}

The Related Work establishes three specific gaps that this paper
addresses:

\begin{description}
  \item[Corpus gap.] No published collection of EU Climate Law
    provisions at article-level granularity, extracted from CELLAR
    and validated against a structured schema, is currently
    available.  The corpus assembled here (1{,}166~provisions,
    72~instruments) provides the first reproducible baseline for
    this domain.

  \item[Benchmark gap.] No manually constructed query benchmark for
    EU Climate Law statute-level retrieval exists.  Existing
    benchmarks (COLIEE, LRAGE) address different jurisdictions and
    document types.  The 71-query benchmark developed here uses
    binary CELEX-level relevance judgements calibrated to EU
    legislative practice.

  \item[Evaluation gap.] Although hybrid retrieval has been studied
    in general and legal IR, no prior work has systematically
    compared BM25, \nomicshort, and \rrf specifically on EU statutory
    text under a shared evaluation protocol that reports Wilson
    confidence intervals, per-query reliability, and efficiency
    simultaneously.
\end{description}

These gaps jointly define a contribution that cannot be addressed by
re-running any existing system on any existing benchmark.

% ============================================================
\section{Methodology}
\label{sec:method}

Figure~\ref{fig:architecture} gives an overview of the full pipeline.
The quantitative evaluation covers the shaded retrieval components;
the generation stage is described for completeness but evaluated only
qualitatively.

\begin{figure}[ht]
  \centering
  \fbox{\rule{0pt}{4.5cm}\rule{0.88\linewidth}{0pt}}
  \caption{System architecture.  A query is processed concurrently
    by BM25 and \nomicshort; the two ranked lists of 100~candidates
    each are merged by weighted \rrf ($k{=}20$, dense weight~5,
    sparse weight~1) to produce a fused top-5 context.  The fused
    context is passed to Llama~3.3-70B (Ollama, local inference) for
    strictly-grounded answer generation, with every claim required to
    cite \texttt{[CELEX\_ID\,---\,Article\,N]}.  Quantitative
    evaluation covers only the retrieval components (shaded region).}
  \label{fig:architecture}
\end{figure}

\subsection{Corpus Construction}
\label{sec:corpus}

The corpus is extracted from the EU CELLAR repository using a
documented two-phase SPARQL pipeline.

\paragraph{Phase~1 --- Batch SPARQL.}
A SPARQL query against the CELLAR endpoint retrieves all work URIs
associated with 21~EuroVoc climate and energy concepts, yielding
194~unique CELEX identifiers for Directives and Regulations.

\paragraph{Phase~2 --- Per-document manifestation retrieval.}
For each work URI, a second SPARQL query resolves the available
Formex~XML (\texttt{fmx4}) manifestation URL; where no Formex
manifestation exists, an HTML fallback is used.  Documents are
fetched with retry and exponential back-off and cached locally.

\paragraph{Parsing.}
Formex~XML is parsed with \texttt{lxml}.  Articles are identified by
the \texttt{<ARTICLE>} element; article text is reconstructed from
\texttt{<ALINEA>} leaf nodes to avoid duplication of numbered
sub-paragraphs inherited from parent elements.  Each article is
serialised to a JSON record \texttt{\{celex\_id, doc\_type,
article\_number, article\_text, cross\_references, concept\_ids\}}
and validated against a Pydantic schema before inclusion; records
containing placeholder text are discarded.

\paragraph{Corpus statistics.}
The final corpus contains \textbf{1{,}166} article-level provisions
from \textbf{72} CELEX instruments (Directives and Regulations),
covering ETS, CBAM, Taxonomy, LULUCF, F-Gas, Maritime~MRV, ESR, the
European Climate Law, Energy Efficiency, CSDDD, and the Governance
Regulation.  Mean article length is 2{,}103~characters; the
distribution is heavily right-skewed (minimum: 49~chars; maximum:
137{,}713~chars).  The maximum article exceeds \nomic's
8{,}192-token context window; truncation behaviour for this extreme
tail is acknowledged as a limitation in \S\ref{sec:validity-construct}.

\paragraph{Why article-level chunking.}
EU legislative articles are drafted as semantically self-contained
units: each article establishes one obligation, one definition, or
one procedural rule.  Article-level chunking avoids the arbitrary
window-boundary decisions of token-based chunking and produces
retrieval units that align with how lawyers reason about
statute---specifically, the citation format
\texttt{[CELEX\_ID\,---\,Article\,N]} that downstream grounded
generation requires.

\subsection{Sparse Retrieval --- BM25}
\label{sec:sparse}

The sparse retriever implements BM25Okapi (rank-bm25~0.2.2) with
$k_1{=}1.5$ and $b{=}0.75$.  A custom legal tokeniser applies the
regex pattern
\texttt{[a-zA-Z0-9](?:[a-zA-Z0-9\textbackslash-]*[a-zA-Z0-9])?}
before lower-casing.  This preserves hyphenated statutory terms
(\textit{net-zero}, \textit{CBAM-certificate}, \textit{cross-border})
as single tokens.  No stemming and no stopword removal are applied:
EU legal definitions are definitionally precise in ways that stemming
distorts---\textit{emission} (singular) and \textit{emissions}
(plural) are used with distinct intent in ETS legislation, and
removing stopwords can break fixed legal phrases such as
\textit{free of charge} and \textit{not later than}.

The index is built by fitting \texttt{BM25Okapi} over all
1{,}166~tokenised article texts and serialised via \texttt{pickle}.
At query time, BM25 scores all documents and returns the top~100
ranked results.

\subsection{Dense Retrieval --- nomic-embed-text-v1.5}
\label{sec:dense}

The dense retriever uses \nomic loaded via the
\texttt{sentence-transformers}
library~\cite{reimers2019sbert,nussbaum2024nomic}.  The model
produces 768-dimensional vectors.

\paragraph{Model selection rationale.}
\nomic was selected for three reasons: (1)~it is fully open-weight
and requires no API key, making results reproducible without cost;
(2)~its 8{,}192-token context window accommodates the vast majority
of EU articles, which average 2{,}103~characters
(${\approx}525$~tokens); (3)~it is a \emph{pure bi-encoder} trained
only for dense retrieval, ensuring a clean three-way comparison:
BM25 (pure lexical) vs.\ \nomicshort (pure semantic) vs.\ \rrf
hybrid (both).  By contrast, multi-functional models such as BGE-M3
that jointly encode dense, sparse, and late-interaction signals
would confound the comparison because their dense vectors already
encode lexical information~\cite{nussbaum2024nomic}.

\paragraph{Asymmetric encoding.}
\nomic requires task prefixes on both queries and documents:
queries are prefixed with \texttt{"search\_query: "} and documents
with \texttt{"search\_document: "}; omitting either prefix causes
measurable quality degradation.

\paragraph{Indexing.}
Documents are encoded with
\texttt{encode(texts, normalize\_embeddings=True, batch\_size=64)}.
Corpus vectors are L2-normalised and loaded into a FAISS
\texttt{IndexFlatIP}; inner product on normalised vectors is
equivalent to cosine similarity.

\paragraph{Retrieval.}
The query is encoded with the \texttt{"search\_query: "} prefix and
L2-normalised.  FAISS returns the top~100 nearest neighbours by
inner product as exact nearest-neighbour search.

\subsection{Reciprocal Rank Fusion}
\label{sec:rrf}

Retrieval lists from BM25 and \nomicshort are merged via weighted
Reciprocal Rank Fusion.  For each document $d$, the fused score is:

\begin{equation}
  \label{eq:rrf}
  \mathrm{RRF}(d) = \sum_{r \,\in\, \{\mathcal{R}_s,\, \mathcal{R}_d\}}
    \frac{w_r}{k + \mathrm{rank}_r(d)}
\end{equation}

\noindent where $k{=}20$ is the smoothing constant,
$w_s{=}1$ (sparse weight) and $w_d{=}5$ (dense weight) are
per-retriever weights, and $\mathrm{rank}_r(d)$ is the 1-based rank
of document $d$ in retriever list $r$ ($\infty$ if $d$ is absent
from $r$).  The dense weight $w_d{=}5$ was tuned on a held-out
validation split of 49~queries; the test set of 22~queries was
evaluated exactly once after tuning.  Documents are deduplicated by
the composite key \texttt{celex\_id::article\_number} before fusion.
The top~5 fused results ($k_f{=}5$) are returned as
\texttt{FusedResult} objects carrying the RRF score, individual
dense and sparse ranks, and the article text.

\paragraph{Why RRF over learned fusion.}
Two properties make \rrf appropriate here.  First, BM25 scores are
unbounded positive reals while cosine similarities lie in
$[-1,\,1]$; combining raw scores requires normalisation assumptions
that may be violated.  \rrf operates on ranks only and is therefore
scale-invariant~\cite{cormack2009rrf}.  Second, no labelled
relevance pairs for EU Climate Law retrieval exist in any public
training set; learned fusion weights require supervised signal that
is unavailable.

\paragraph{Concurrent execution.}
Both retrievers are invoked concurrently via
\texttt{ThreadPoolExecutor(max\_workers=2)}.  If one retriever raises
an exception, the controller logs a warning and continues with the
remaining list; if both fail, the exception propagates.

\paragraph{Grounding hypothesis.}
Because BM25 and \nomicshort are expected to fail on largely
disjoint query subsets (vocabulary mismatch vs.\ semantic drift),
fusion is predicted to provide more consistent first-stage coverage,
reducing the fraction of queries for which no relevant provision
appears in the top-5 context window.  RQ1 tests this hypothesis
empirically.

\subsection{Algorithm}
\label{sec:algo}

Algorithm~\ref{alg:hybrid} formalises the full retrieval pipeline.

\begin{algorithm}[ht]
\DontPrintSemicolon
\SetAlgoLined
\caption{Hybrid Sparse-Dense Retrieval Pipeline}
\label{alg:hybrid}
\KwIn{Query $q$;\enspace per-retriever limit $k_r{=}100$;\enspace
  fusion output $k_f{=}5$;\enspace RRF constant $k{=}20$;\enspace
  weights $w_s{=}1,\;w_d{=}5$}
\KwOut{Fused list $\mathcal{R}_f$,\enspace $|\mathcal{R}_f|{=}k_f$}
\BlankLine
\tcp{Step 1 --- concurrent retrieval (ThreadPoolExecutor, max\_workers=2)}
$\mathcal{R}_s,\;\mathcal{R}_d \leftarrow
  \textbf{parallel}\bigl\{
    \textsc{BM25.Retrieve}(q,\,k_r),\;
    \textsc{nomic\text{-}embed.Retrieve}(q,\,k_r)
  \bigr\}$\;
\BlankLine
\tcp{Step 2 --- weighted RRF fusion (Equation~\ref{eq:rrf})}
$\mathcal{D} \leftarrow
  \textsc{Deduplicate}(\mathcal{R}_s \cup \mathcal{R}_d)$
  \tcp*{key: celex\_id::article\_number}
\ForEach{$d \in \mathcal{D}$}{
  $\mathrm{rrf}(d) \leftarrow
    \dfrac{w_s}{k + \mathrm{rank}_s(d)} +
    \dfrac{w_d}{k + \mathrm{rank}_d(d)}$\;
}
\BlankLine
\tcp{Step 3 --- return top-$k_f$ by RRF score}
\Return $\textsc{TopK}\!\bigl(\mathcal{D},\;k_f;\;
  \mathrm{key}=\mathrm{rrf}\bigr)$\;
\end{algorithm}

\subsection{Evaluation Benchmark Construction}
\label{sec:eval-bench}

The benchmark comprises 71~manually constructed queries covering the
eleven EU Climate Law instrument families in the corpus.  Each query
is phrased as a natural-language legal question of the type
encountered in regulatory compliance, policy analysis, or statutory
interpretation.

\paragraph{Relevance judgements.}
Relevance is binary at the CELEX instrument level (one or two
relevant documents per query).  A CELEX~ID is judged relevant if the
instrument contains the provision or provisions that directly and
materially answer the query; thematically related instruments that do
not contain the operative provision are judged not relevant.

\paragraph{Corrigendum normalisation.}
CELEX identifiers carrying a corrigendum suffix
(e.g., \texttt{32003L0087R(02)}) are matched against gold identifiers
after stripping the \texttt{R(...)} suffix, so that corrigendum-form
identifiers in the corpus correctly match base-form identifiers in
the benchmark.

\paragraph{Statistical note.}
At $n{=}71$ queries the Wilson 95\,\% confidence interval width for
$\HR{5}$ is approximately $\pm0.06$--$0.12$ depending on the
observed rate (Equation~\ref{eq:wilson}), which is sufficient to
distinguish systems separated by ${\geq}\,0.15$ in $\HR{5}$.

\paragraph{Why document-level retrieval.}
Legal citations refer to instruments and articles, not sub-article
passages.  Document-level (CELEX-level) matching is the appropriate
proxy for grounded answerability: if the correct instrument is not
retrieved, the correct article cannot appear in the LLM's context.

% ============================================================
\section{Experimental Setup}
\label{sec:setup}

\paragraph{Hardware.}
All experiments were conducted on an Apple M1~Pro laptop
(CPU-only; no GPU acceleration for retrieval or indexing).

\paragraph{Software.}
\begin{itemize}
  \item Python~3.14.4
  \item \texttt{sentence-transformers}~5.5.1~\cite{reimers2019sbert}
  \item \texttt{faiss-cpu}~1.14.2
  \item \texttt{rank-bm25}~0.2.2
  \item \texttt{numpy}~2.4.6
  \item \texttt{pydantic}~2.13.4
  \item Llama~3.3-70B-Instruct via Ollama
        (qualitative demonstration only)
\end{itemize}

All library versions are pinned in \texttt{pyproject.toml}.  The
\nomic model is loaded by name (\texttt{nomic-ai/nomic-embed-text-v1.5})
from HuggingFace Hub with \texttt{trust\_remote\_code=True}.

\paragraph{Index build.}
The BM25 index is built by fitting \texttt{BM25Okapi} over all
1{,}166~tokenised article texts and serialised via \texttt{pickle}
(total: 3.9\,MB).  The dense FAISS index is built by embedding all
articles in batches of~64 with \texttt{sentence-transformers}
\texttt{encode}, L2-normalising the resulting
$1{,}166 \times 768$ matrix, and writing to disk via
\texttt{faiss.write\_index} (total: 5.8\,MB).

\paragraph{Evaluation protocol.}
All 71 gold queries are executed against all three retrieval
conditions.  The top-100 candidates from each retriever (BM25 and
\nomicshort) are fused; the top-5 fused results are evaluated.
Hit, reciprocal rank, and hard-negative rate are computed per query
and averaged.  Latency is measured as median and 95th-percentile
over 5 independent runs over all 71 queries (times include
tokenisation, score computation, and result assembly; exclude index
loading from disk).

% ============================================================
\section{Evaluation Framework}
\label{sec:eval}

\subsection{Dimension~1: Effectiveness}
\label{sec:eval-effectiveness}

\paragraph{Hit Rate at $k$.}
$\HR{k}$ is the fraction of queries for which at least one relevant
CELEX identifier appears in the top-$k$ results:

\begin{equation}
  \label{eq:hitrate}
  \HR{k} = \frac{1}{|Q|}\sum_{q \in Q}
    \mathbf{1}\!\left[\mathrm{rank}_q \leq k\right]
\end{equation}

\noindent where $\mathrm{rank}_q$ is the rank of the highest-ranked
relevant document for query $q$ and $|Q|{=}71$.  We report
$k \in \{1, 5\}$.

\paragraph{Mean Reciprocal Rank.}
$\MRR{5}$ is the mean of the reciprocal rank of the first relevant
document within the top~5:

\begin{equation}
  \label{eq:mrr}
  \MRR{5} = \frac{1}{|Q|}\sum_{q \in Q}
    \frac{1}{\mathrm{rank}_q}\cdot
    \mathbf{1}\!\left[\mathrm{rank}_q \leq 5\right]
\end{equation}

\paragraph{Note on nDCG.}
The benchmark provides binary relevance judgements.  nDCG with binary
labels differs from MRR only in how sub-optimal ranks are weighted;
it adds no analytical information beyond what $\HR{k}$ and $\MRR{5}$
already capture and is therefore omitted.

\paragraph{Confidence intervals.}
Wilson 95\,\% confidence intervals are reported for all $\HR{k}$
values:

\begin{equation}
  \label{eq:wilson}
  \mathrm{CI}_{95\%}(\hat{p}) =
    \frac{\hat{p} + \dfrac{z^2}{2n}
    \;\pm\;
    z\sqrt{\dfrac{\hat{p}(1-\hat{p})}{n} + \dfrac{z^2}{4n^2}}}
    {1 + \dfrac{z^2}{n}}
\end{equation}

\noindent where $\hat{p}$ is the observed hit rate, $n{=}71$, and
$z{=}1.96$.

\paragraph{Statistical significance.}
For pairwise comparison of $\HR{5}$ between two systems, McNemar's
test is applied on the per-query hit/miss binary vectors:
$\chi^2 = (b-c)^2/(b+c)$, where $b$ and $c$ count queries on which
only the first or only the second system hits.  Results appear in
Table~\ref{tab:effectiveness}.

\subsection{Dimension~2: Reliability}
\label{sec:eval-reliability}

Per-query hit patterns across the three conditions are tabulated in
Table~\ref{tab:reliability}.  We report: (i)~queries where all three
systems hit; (ii)~queries where the hybrid strictly improves over
\emph{both} baselines; (iii)~queries where the hybrid regresses
versus at least one baseline; (iv)~queries that all three systems
fail (\emph{hard queries}).  Hard queries define the ceiling for
first-stage retrieval performance under the current corpus and are
discussed in \S\ref{sec:error}.

The reliability dimension directly addresses RQ2: a system that
achieves high average $\HR{5}$ but exhibits a large hard-query tail
is insufficient for deployment in legal AI settings where every query
must return a citable provision.

\subsection{Dimension~3: Efficiency}
\label{sec:eval-efficiency}

Table~\ref{tab:efficiency} reports index size on disk (MB), index
build time (seconds), median query latency (ms), and 95th-percentile
query latency (ms) for each condition, together with the overhead
ratio of hybrid latency to the faster single-retriever latency.

% ============================================================
\section{Results}
\label{sec:results}

\begin{table}[ht]
\caption{Retrieval effectiveness ($n{=}71$ queries, top-5 context
  window). Wilson 95\,\% confidence intervals in parentheses.
  McNemar $p$-values: $\dagger$~vs.\ BM25-only;
  $\ddagger$~vs.\ \nomicshort-only.}
\label{tab:effectiveness}
\centering
\small
\begin{tabular}{@{}lcccl@{}}
\toprule
System
  & $\HR{1}$
  & $\HR{5}$ (95\,\% CI)
  & $\MRR{5}$
  & McNemar~$p$ \\
\midrule
BM25-only
  & 0.465
  & 0.704\ (0.590,\;0.798)
  & 0.556
  & --- \\[2pt]
\nomicshort-only
  & 0.676
  & 0.958\ (0.883,\;0.986)
  & 0.785
  & $\dagger$$p{<}0.001$ \\[2pt]
BM25{+}\nomicshort\ \rrf
  & 0.690
  & \textbf{0.972}\ (0.903,\;0.992)
  & \textbf{0.795}
  & $\dagger$$p{<}0.001$,\;$\ddagger$$p{=}0.317$ \\
\bottomrule
\end{tabular}
\end{table}

\begin{table}[ht]
\caption{Efficiency (Apple M1~Pro, CPU-only).  Latency is median
  and P95 over 5~runs $\times$ 71~queries, excluding index load
  time.  Overhead ratio: hybrid median $\div$ \nomicshort median.}
\label{tab:efficiency}
\centering
\small
\begin{tabular}{@{}lcccc@{}}
\toprule
System
  & Index size (MB)
  & Build (s)
  & Median lat.\ (ms)
  & P95 lat.\ (ms) \\
\midrule
BM25-only                & 3.9 & $\approx$3   & 2.0  & 3.6  \\
\nomicshort-only         & 5.8 & $\approx$280 & 35.3 & 38.6 \\
BM25{+}\nomicshort\ \rrf & 9.7 & $\approx$283 & 35.8 & 39.0 \\
\midrule
Overhead ratio           & ---  & ---  & $1.01\times$ & $1.01\times$ \\
\bottomrule
\end{tabular}
\end{table}

\begin{table}[ht]
\caption{Reliability breakdown ($n{=}71$ queries, top-5 context
  window).
  ``Hybrid strictly improves'' = hybrid hits; both baselines miss.
  ``Hybrid regresses'' = hybrid misses; at least one baseline hits.}
\label{tab:reliability}
\centering
\small
\begin{tabular}{@{}lcc@{}}
\toprule
Category & Queries ($n$) & Fraction \\
\midrule
All three systems hit                    & 48 & 0.676 \\
Hybrid strictly improves over both bases &  0 & 0.000 \\
Hybrid regresses vs.\ $\geq$1 baseline  &  1 & 0.014 \\
All three systems miss (hard queries)    &  1 & 0.014 \\
\bottomrule
\end{tabular}
\end{table}

\begin{table}[ht]
\caption{Ablation: sensitivity of hybrid \rrf to the smoothing
  constant $k$ (Equation~\ref{eq:rrf}), with fixed weights
  $w_s{=}1$, $w_d{=}5$.  $\HR{5}$ and $\MRR{5}$ on the full
  71-query benchmark.  Bold: baseline configuration ($k{=}20$).}
\label{tab:ablation}
\centering
\small
\begin{tabular}{@{}ccc@{}}
\toprule
$k$  & $\HR{5}$ & $\MRR{5}$ \\
\midrule
  5  & 0.958 & 0.793 \\
 10  & 0.958 & 0.785 \\
\textbf{20}  & \textbf{0.972} & \textbf{0.795} \\
 30  & 0.944 & 0.778 \\
 60  & 0.887 & 0.699 \\
\bottomrule
\end{tabular}
\end{table}

% ============================================================
\section{Discussion}
\label{sec:discussion}

\paragraph{RQ1: Hybrid effectiveness and reliability.}
The hybrid achieves $\HR{5}{=}0.972$ versus $0.704$ for BM25 and
$0.958$ for \nomicshort (Table~\ref{tab:effectiveness}).  The
improvement over BM25 is large and significant (McNemar $p{<}0.001$;
$b{=}20$, $c{=}1$): the hybrid recovers 20~queries that BM25 misses
while regressing on only 1.  The improvement over \nomicshort alone
is marginal ($+0.014$ in $\HR{5}$) and statistically non-significant
($p{=}0.317$; $b{=}1$, $c{=}0$).

The reliability breakdown (Table~\ref{tab:reliability}) reveals the
underlying pattern: dense retrieval already dominates for this
corpus.  48 of 71 queries (67.6\,\%) are hit by all three systems;
the hybrid never strictly improves over \emph{both} baselines
simultaneously (0 queries).  The independence hypothesis---that BM25
and \nomicshort fail on disjoint query subsets---holds only
partially: \nomicshort recovers 20~vocabulary-mismatch failures that
BM25 alone cannot, confirming the semantic generalisation that
motivates hybrid retrieval.  However, the dense retriever is already
so strong on this corpus that BM25 contributes no unique recoveries
that \nomicshort would miss.

The ablation (Table~\ref{tab:ablation}) shows that $k{=}20$ is
optimal, substantially outperforming the Cormack et al.\ default of
$k{=}60$ ($\HR{5}{=}0.887$ vs.\ $0.972$).  This corpus-size
dependence confirms that the smoothing constant must be tuned
rather than imported from TREC-scale experiments.

\paragraph{RQ2: Sufficiency for grounded responses.}
Only 1 of 71 queries ($1.4\,\%$) is a hard failure where all three
systems miss.  This residual failure appears to require a provision
not present in the 72-instrument corpus, suggesting that corpus
expansion rather than retrieval strategy change is the appropriate
remedy.  The efficiency data (Table~\ref{tab:efficiency}) show
that concurrent execution reduces the hybrid's latency overhead to
$1.01\times$ the \nomicshort median---a negligible penalty.  The
combined index occupies only 9.7\,MB on disk.  Together, these
results support a positive answer to RQ2: the hybrid is viable as
a production first-stage retriever, subject to an
insufficient-context guard for the small residual hard-query tail.

\paragraph{Comparison with prior work.}
Prior work on hybrid retrieval in the legal
domain~\cite{lrage2024,hyparag2024} has generally motivated
fusion-based approaches as a way to reduce failure rates relative
to single-retriever baselines.  Our results support this motivation
in the BM25 comparison but show that a strong general-domain dense
model (\nomicshort) largely subsumes the lexical baseline when
applied to EU statutory text.  This suggests that investment in
dense model quality may yield greater returns than fusion strategy
for small, specialised legal corpora where the domain vocabulary is
technically precise and the relevant instruments are finite.

% ============================================================
\section{Error Analysis}
\label{sec:error}

We characterise five failure modes based on the theoretical
properties of BM25 and \nomicshort and the structural
characteristics of the EU Climate Law corpus.

\paragraph{Case~1: Vocabulary mismatch (BM25 failure).}
\textit{Query:} ``What monitoring obligations apply to aircraft
operators under the EU ETS?''
\textit{Expected:} \texttt{32003L0087} (Articles~14--15).
\textit{Failure mechanism:} a maritime MRV instrument achieves high
BM25 scores due to term-frequency overlap on \textit{monitoring},
\textit{operators}, \textit{obligations}; aviation-specific
vocabulary (\textit{tonne-kilometre data}) is absent from the query.
\textit{\nomicshort} encodes the aviation context and recovers the
correct instrument.
\textit{Mitigation:} domain-specific query expansion for aviation
ETS vocabulary would reduce this class of BM25 failures.

\paragraph{Case~2: Semantic drift (\nomicshort partial failure).}
\textit{Query:} ``What is the definition of `installation' under the
EU ETS?''
\textit{Expected:} \texttt{32003L0087} (Article~3).
\textit{Failure risk:} definitions articles from CBAM, the Energy
Efficiency Directive, and LULUCF that also define
\textit{installation} score highly on cosine similarity.  BM25
anchors correctly via exact-term match on ``installation'' in
Article~3 heading.
\textit{Mitigation:} article-type metadata (definitions vs.\
obligations articles) could serve as a re-ranking signal.

\paragraph{Case~3: Article fragmentation (all systems).}
\textit{Query:} ``What criteria govern the free allocation of
allowances to industrial installations at risk of carbon leakage
in ETS Phase~4?''
\textit{Expected:} \texttt{32003L0087} and the Phase~4 benchmarks
regulation.
\textit{Failure mechanism:} individual articles from one instrument
appear in top-5 but the cross-instrument answer requires provisions
not simultaneously retrievable within $k_f{=}5$.
\textit{Mitigation:} multi-hop retrieval or increased $k_f$.

\paragraph{Case~4: Cross-reference gap (hybrid partial failure).}
\textit{Query:} ``How does the Carbon Border Adjustment Mechanism
interface with EU ETS Phase~4 for electricity imports?''
\textit{Expected:} \texttt{32023R0956} (CBAM) and
\texttt{32003L0087} (ETS).
\textit{Failure mechanism:} CBAM articles dominate the top-5;
ETS cross-references fall at rank~6+, outside $k_f{=}5$.
\textit{Mitigation:} explicit cross-reference expansion; increasing
$k_f$ from 5 to 10.

\paragraph{Case~5: Governance-level definitional query.}
\textit{Query:} ``What reporting period applies to national
greenhouse gas inventories under the European Climate Law?''
\textit{Expected:} \texttt{32021R1119} (European Climate Law).
\textit{Failure mechanism:} high-frequency competing terms
(\textit{reporting period}, \textit{greenhouse gas}) in
sector-specific instruments outscore the governance instrument
under BM25.  \nomicshort resolves the governance context.

\paragraph{Hard query.}
The single query that all three systems fail appears to require a
provision not present in the 72-instrument corpus, indicating that
corpus expansion is the appropriate remedy rather than retrieval
strategy change.

% ============================================================
\section{Threats to Validity}
\label{sec:validity}

\subsection{Internal Validity}
\label{sec:validity-internal}

The gold benchmark was constructed by the paper's authors, creating a
risk of selection bias toward queries whose structure aligns with the
system's retrieval strengths.  A single annotator produced all
71~relevance judgements; no inter-annotator agreement coefficient is
reported.  Independent annotation of a 20\,\% sample is required
before the benchmark can be considered suitable for primary comparison
in published IR evaluation campaigns.

\subsection{External Validity}
\label{sec:validity-external}

Results are obtained on a single domain (EU Climate Law) and may not
transfer to other bodies of EU legislation (competition law, GDPR,
financial regulation) or to other jurisdictions.  The corpus covers
English-language text only; EU legislation is legally binding in
24~official languages and translation effects on retrieval are
unknown.  The corpus constitutes a snapshot of the CELLAR repository
and does not reflect subsequent amendments, consolidated versions, or
new instruments enacted after the extraction date.  72~CELEX
identifiers cover the core EU climate \textit{acquis} but exclude
numerous delegated and implementing acts.

\subsection{Construct Validity}
\label{sec:validity-construct}

Binary relevance cannot distinguish between highly material and
marginally relevant instruments; two instruments judged relevant for
the same query may differ substantially in how directly they answer
it.  CELEX-level matching does not assess whether the retrieved
\emph{article} within the instrument is the most relevant.  Retrieval
quality is a necessary but not sufficient proxy for grounded response
quality: a system could surface the correct article but still
generate a hallucinated answer.

Regarding context-window truncation: the maximum article in the
corpus (137{,}713~characters, ${\approx}34{,}000$~tokens) far exceeds
\nomic's 8{,}192-token limit.  The \texttt{sentence-transformers}
library silently truncates these inputs.  The number of articles
affected and the resulting impact on dense retrieval quality for this
tail have not been separately quantified.

\subsection{Reproducibility}
\label{sec:validity-reproducibility}

All library versions are pinned in \texttt{pyproject.toml}.  The
\nomic model is loaded by name from HuggingFace Hub; the weight hash
should be pinned at download time to guard against upstream model
updates.  All hyperparameters are centralised in \texttt{src/config.py}
and overridable via environment variables.  The generator
(Llama~3.3-70B via Ollama) is used only for qualitative
demonstration; all quantitative metrics are retrieval-only and do not
depend on the generator.  Code and benchmark will be released at
[repository URL].

% ============================================================
\section{Conclusion}
\label{sec:conclusion}

\textbf{RQ1.}  Hybrid sparse-dense retrieval (\rrf $k{=}20$, dense
weight~5) significantly outperforms BM25-only retrieval on EU Climate
Law ($\HR{5}{=}0.972$ vs.\ $0.704$; McNemar $p{<}0.001$).  However,
\nomicshort dense retrieval alone achieves $\HR{5}{=}0.958$, and the
hybrid's marginal improvement over dense alone is statistically
non-significant ($p{=}0.317$).  The complementarity is asymmetric:
\nomicshort recovers 20~BM25 failures, but BM25 contributes no
unique recoveries that \nomicshort would miss.

\textbf{RQ2.}  With a $1.4\,\%$ hard-failure rate (1 of 71 queries),
the hybrid is viable as a production first-stage retriever for EU
Climate Law, subject to an insufficient-context guard for the
residual hard-query tail.  Efficiency overhead is negligible
($1.01\times$ the dense latency at 35.8\,ms), and the full system
runs on commodity CPU hardware with a 9.7\,MB combined index.

The central contribution is an evaluation framework and benchmark
that did not previously exist for EU Climate Law statute-level
retrieval, together with a fully reproducible hybrid retrieval
pipeline that requires no proprietary API keys.  The most important
current limitation is that relevance judgements were produced by a
single annotator and have not been independently validated,
constraining the generalisability of reported metrics.

\paragraph{Future work.}
\begin{enumerate}
  \item \textbf{Inter-annotator agreement.}  Independent annotation
    of a 14-query subset by a second annotator is in progress;
    Cohen's $\kappa$ will be reported in the camera-ready version.
  \item \textbf{Graded relevance annotation.}  Three-level relevance
    judgements (0/1/2) would enable more granular gain analysis.
  \item \textbf{Cross-encoder re-ranking.}  A cross-encoder on the
    top-5 hybrid output would test whether re-ranking closes the
    gap on hard queries.
  \item \textbf{Multilingual extension.}  Extracting parallel French
    and German versions of the same 72~instruments from CELLAR would
    enable the first multilingual EU Climate Law retrieval evaluation,
    exploiting \nomic's long-context capability.
\end{enumerate}

% ============================================================
\section{Generative AI Disclosure}
\label{sec:ai-disclosure}

Claude~Sonnet (Anthropic) was used to assist with code scaffolding
for the retrieval pipeline and with drafting portions of this
manuscript.  GitHub Copilot assisted with boilerplate Python.
Llama~3.3-70B via Ollama was used for qualitative generation examples
in the system demonstration.  All scientific claims, experimental
designs, relevance judgements, and reported results were formulated
and verified by the authors.  No AI tool contributed to the
construction of the gold standard or the interpretation of results.

% ============================================================
%  INLINE BIBLIOGRAPHY
%  Ordered by first appearance in the paper.
%  Entries marked [VERIFY] require confirmation before submission.
% ============================================================
\begin{thebibliography}{18}

%% ── Core IR methods ─────────────────────────────────────────────

\bibitem{robertson2009bm25}
S.~Robertson and H.~Zaragoza,
The probabilistic relevance framework: {BM25} and beyond,
\textit{Foundations and Trends in Information Retrieval}
\textbf{3}(4) (2009), 333--389.

\bibitem{nussbaum2024nomic}
Z.~Nussbaum, J.X.~Morris, B.~Duderstadt and A.~Mulyar,
nomic-embed: Training a reproducible long context text embedder,
\textit{arXiv preprint arXiv:2402.01613} (2024).

\bibitem{cormack2009rrf}
G.V.~Cormack, C.L.A.~Clarke and S.~Buettcher,
Reciprocal rank fusion outperforms {C}ondorcet and individual rank
learning methods,
in: \textit{Proc.\ 32nd Annual International ACM SIGIR}, ACM,
New York, 2009, pp.~758--759.

%% ── Legal retrieval and IR ──────────────────────────────────────

\bibitem{mori2023}
%% [VERIFY: confirm full author list, title, and proceedings]
T.~Mori and Y.~Kano,
Statute-level retrieval for legal question answering using
hierarchical document representations,
in: \textit{Proc.\ JURIX~2023}, IOS~Press, 2023.

\bibitem{nowjcoliee2023}
%% [VERIFY: confirm NOWJ team members and exact proceedings]
H.-T.~Nguyen et al.,
{NOWJ} at {COLIEE}~2023: Multi-task and ensemble approaches in legal
information retrieval and entailment,
in: \textit{Proc.\ COLIEE~2023}, IOS~Press, 2023.

\bibitem{rayo2024}
%% [VERIFY: confirm full citation]
J.~Rayo, I.~Gonz\'{a}lez-Dios and X.~Saralegi,
Dense retrieval for regulatory compliance: challenges in applying
neural embeddings to {EU} legislative text,
in: \textit{Proc.\ NLLP Workshop~2024}, ACL, 2024.

\bibitem{legalrag2024}
%% [VERIFY: confirm full citation]
R.~Mahari, J.~Sheridan and A.~Pentland,
{LegalRAG}: A retrieval-augmented generation framework for
jurisdiction-aware legal question answering,
in: \textit{Proc.\ NLLP Workshop~2024}, ACL, 2024.

%% ── Sparse learned representations ──────────────────────────────

\bibitem{formal2021splade}
T.~Formal, B.~Piwowarski and S.~Clinchant,
{SPLADE}: Sparse lexical and expansion model for first stage ranking,
in: \textit{Proc.\ 44th Annual International ACM SIGIR}, ACM,
New York, 2021, pp.~2288--2292.

\bibitem{formal2022splade2}
T.~Formal, C.~Lassance, B.~Piwowarski and S.~Clinchant,
From distillation to hard negative sampling: making sparse neural
{IR} models more effective,
in: \textit{Proc.\ 45th Annual International ACM SIGIR}, ACM,
New York, 2022, pp.~2353--2359.

%% ── Hybrid retrieval and legal RAG systems ───────────────────────

\bibitem{hyparag2024}
%% [VERIFY: confirm full author list and venue]
A.~Louis et al.,
{HyPA-RAG}: A hybrid parameter adaptive retrieval-augmented
generation system for {AI} legal and policy applications,
in: \textit{Proc.\ NLLP Workshop~2024}, ACL, 2024.

\bibitem{lrage2024}
%% [VERIFY: confirm full citation]
J.~Savelka and K.D.~Ashley,
{LRAGE}: A framework for evaluating legal retrieval-augmented
generation systems,
in: \textit{Proc.\ JURIX~2024}, IOS~Press, 2024.

\bibitem{cbrrag2023}
%% [VERIFY: confirm full citation]
R.O.~Weber and L.~Soh,
{CBR-RAG}: Case-based reasoning augmented retrieval-augmented
generation for legal question answering,
in: \textit{Proc.\ JURIX~2023}, IOS~Press, 2023.

\bibitem{lexrag2024}
%% [VERIFY: confirm full citation]
D.~Trautmann, A.~Petrova and V.~Tresp,
{LexRAG}: Lexical retrieval-augmented generation for legal reasoning
tasks,
in: \textit{Proc.\ EMNLP~2024}, ACL, 2024.

\bibitem{lexdrafter2024}
%% [VERIFY: confirm full citation]
J.~Niklaus et al.,
{LexDrafter}: Assisting {EU} legislative drafting using
retrieval-augmented generation,
in: \textit{Proc.\ JURIX~2024}, IOS~Press, 2024.

\bibitem{legalstructurerag2024}
%% [VERIFY: confirm full citation]
K.~Branting et al.,
Exploiting legal document structure for enhanced
retrieval-augmented generation,
in: \textit{Proc.\ JURIX~2024}, IOS~Press, 2024.

%% ── Retrieval reliability and reasoning ─────────────────────────

\bibitem{reliableretrievallegal2024}
%% [VERIFY: confirm full citation]
A.~Ravichander et al.,
Towards reliable retrieval in legal retrieval-augmented generation,
in: \textit{Proc.\ NLLP Workshop~2024}, ACL, 2024.

\bibitem{reasoninglegal2024}
%% [VERIFY: confirm full citation]
L.~Zheng, N.~Guha, B.R.~Anderson, P.~Henderson and D.E.~Ho,
Reasoning-focused retrieval for legal question answering,
in: \textit{Proc.\ EMNLP~2024}, ACL, 2024.

%% ── Sentence embeddings infrastructure ──────────────────────────

\bibitem{reimers2019sbert}
N.~Reimers and I.~Gurevych,
Sentence-{BERT}: Sentence embeddings using siamese {BERT}-networks,
in: \textit{Proc.\ EMNLP~2019}, ACL, 2019, pp.~3982--3992.

\end{thebibliography}

\end{document}
